\section{Preliminaries}
\label{sec:preliminaries}

\subsection{Notation}

We work over a finite field $\mathbb{F}$ (KoalaBear in our deployment,
$|\mathbb{F}| \approx 2^{31}$) and an extension $\mathbb{F}_q$ with
$q = |\mathbb{F}|^D$ for $D=4$, giving challenger soundness
$\sim 2^{-124}$.  $N = 2^n$ denotes a power-of-two domain size,
typically the height of a chip's trace.  Bold lowercase letters
($\mathbf{x}, \mathbf{r}, \mathbf{z}$) denote vectors of field
elements; uppercase ($P, Q, F$) denote polynomials.

For a vector of evaluations $\mathbf{v} \in \mathbb{F}^{2^m}$, the
\emph{multilinear extension} (MLE) $\widetilde{\mathbf{v}}:
\mathbb{F}^m \to \mathbb{F}$ is the unique multilinear polynomial
satisfying $\widetilde{\mathbf{v}}(\mathbf{b}) = v_{\mathrm{idx}
(\mathbf{b})}$ for $\mathbf{b} \in \{0,1\}^m$.  We use the
multilinear equality polynomial $\mathsf{eq}(\mathbf{r}, \mathbf{x})
= \prod_{i=1}^{m} (r_i x_i + (1-r_i)(1-x_i))$, which satisfies
$\widetilde{\mathbf{v}}(\mathbf{r}) = \sum_{\mathbf{b} \in \{0,1\}^m}
\mathsf{eq}(\mathbf{r}, \mathbf{b}) \cdot v_{\mathrm{idx}(\mathbf{b})}$.

\subsection{STARKs and the Ziren Pipeline}

A STARK proves correct execution of a virtual-machine program by:
(1) expressing the program's per-cycle behaviour as an algebraic
intermediate representation (AIR), (2) committing to the resulting
trace via a polynomial commitment scheme (PCS), (3) verifying that
the committed trace satisfies the AIR's transition and boundary
constraints (via a quotient polynomial argument or a sumcheck),
and (4) verifying lookups (via the LogUp~\cite{logup-gkr} family of
arguments).

Ziren's main branch implements step~(2) via FRI on a multiplicative
subgroup, with separate commits to main, permutation, and quotient
traces.  Step~(3) is the standard quotient identity
$Q(\zeta) \cdot Z_H(\zeta) = C_{\mathrm{batched}}(\zeta)$ at a
random univariate point $\zeta$.

\subsection{Sumcheck Protocol}

The sumcheck protocol~\cite{sumcheck} reduces a claim about a
polynomial sum over the Boolean hypercube to a single polynomial
evaluation:
\[
  \text{Claim:}\quad \sum_{\mathbf{b} \in \{0,1\}^m} F(\mathbf{b}) = v.
\]
For each round $i \in [m]$, the prover sends a univariate polynomial
$p_i(X) = \sum_{\mathbf{b}' \in \{0,1\}^{m-i}}
F(r_1, \ldots, r_{i-1}, X, \mathbf{b}')$ of degree $\deg(F)$, and
the verifier checks $p_i(0) + p_i(1) = s_i$ (where $s_1 = v$ and
$s_{i+1} = p_i(r_i)$ for verifier-sampled $r_i \in \mathbb{F}_q$).
After $m$ rounds, the verifier checks
$F(r_1, \ldots, r_m) = s_{m+1}$.  Soundness error is
$d \cdot m / |\mathbb{F}_q|$ for $\deg(F) = d$.

\subsection{WHIR}

WHIR~\cite{whir2024} is an interactive oracle proof of proximity
(IOPP) for Reed--Solomon (RS) codes that supports multilinear
polynomial commitments with $O(\log n)$ verifier work and proof
sizes $5$--$10\times$ smaller than FRI at equivalent security.

A WHIR commitment to $f \in \mathbb{F}^{2^n}$ encodes $f$'s
coefficients into an RS codeword on a coset of size $2^{n+r}$
($r$ = log-inv-rate) and merkle-commits the row hashes.  The open
phase combines sumcheck folding rounds with STIR query openings.

The Plonky3 implementation \texttt{p3-whir} provides a
\texttt{MultilinearPcs} adapter requiring opening points
\emph{at commit time} (absorbed into Fiat-Shamir before the
batching challenge $\alpha$ is sampled).  This is the standard
soundness ordering when opening points are pre-known; our
work (\S\ref{sec:whir-late-binding}) constructs an alternative
\emph{late-binding} ordering for protocols where the opening
point is derived post-commit.

\subsection{LogUp-GKR}

LogUp~\cite{logup-gkr} reduces a multiset-equality lookup argument
to a sum-of-fractions identity
\[
  \sum_{i \in \mathrm{senders}} \frac{m_i}{\alpha - f_i} \;=\;
  \sum_{j \in \mathrm{receivers}} \frac{m_j}{\alpha - f_j},
\]
where $m_*$ are multiplicities and $f_*$ are interaction
fingerprints.  GKR~\cite{logup-gkr} verifies this fraction-sum
identity via a layered tree of fraction combinations
$(a, b) \oplus (c, d) = (ad + bc, bd)$, opening only the root
$(\mathrm{num}, \mathrm{denom})$ and verifying $\mathrm{num} = 0$.

The verification reduces tree-root claims to leaf claims via
sumcheck rounds at each layer.  The leaf claims are MLE evaluations
of multiplicities and fingerprints at a sumcheck-derived point
$r_{\mathrm{row}}$, which is therefore \emph{post-commit}.  This
post-commit derivation is the central structural difficulty for
binding the leaf claims to a WHIR commitment of the trace.

\subsection{Plonky3 Architecture}

Ziren is built on Plonky3~\cite{plonky3}, which provides chip-based
AIR composition, Fiat-Shamir challenger, MMCS (mixed-matrix
commitment scheme) merkle trees, and a recursion compiler.  Our
work modifies neither the chip composition nor the recursion
compiler core; we add new chips (sumcheck verification),
new instructions (sumcheck verify), new prover/verifier paths
(WHIR fast path), and new helpers (late-binding, cross-binding,
jagged + late-binding) on top of the existing infrastructure.
